% generated by bibm/supplement/make_sm_tables.py -- do not edit; sources: gue_family_summary.csv
\begin{table}[!htbp]
\caption{GUE task families: how the negatives were built, how much of the DNA is shared with GB/NT, and the labels. Source: \src{gue\_family\_summary.csv}.}
\label{tab:guefam}
\centering
\scriptsize
\setlength{\tabcolsep}{3pt}
\begin{tabular}{>{\raggedright\arraybackslash}p{0.13\linewidth}r>{\raggedright\arraybackslash}p{0.33\linewidth}l>{\raggedright\arraybackslash}p{0.17\linewidth}rr}
\toprule
Family & $n$ & Negatives / label-0 class (as documented) & Shared & Labels & $k$-mer & DNABERT-2 \\
\midrule
Core promoter (70 bp) & 3 & as prom\_300, center -34..+35 bp around the TSS & 50.1--77.4\% & hard:2, positional:1 & 58.2 & 70.5 \\
Promoter (300 bp) & 3 & non-TATA: randomly substituted promoter sequences; TATA: random non-promoter genomic sequences carrying a TATA motif & 50.0--65.8\% & solved:2, positional:1 & 66.8 & 84.2 \\
TF binding, human & 5 & non-overlapping genomic sequences of the same length and GC content (GC-matched) & 19.8--44.8\% & hard:3, solved:2 & 63.1 & 70.1 \\
Splice site (3-class) & 1 & Wang et al. 2019 donor/acceptor/non-site, reconstructed with adversarial negatives & 68.4--68.4\% & positional:1 & 34.7 & 85.0 \\
TF binding, mouse & 5 & dinucleotide-shuffled positives (dinucleotide frequencies preserved) & 0.0--0.2\% & hard:3, solved:2 & 63.5 & 68.0 \\
Yeast histone marks & 10 & Pham et al. 2005 datasets (jaist.ac.jp/\textasciitilde{}tran/nucleosome): yeast genomic windows (ORF and intergenic ChIP-chip regions) with high (label 1) vs low (label 0) relative histone occupancy/modification; no sampled negatives; GUE split 8:1:1 & 100.0--100.0\% & shared-motif:5, contamination:4, solved:1 & 51.5 & 56.0 \\
Covid variant (9-class) & 1 & 9 SARS-CoV-2 variants (GISAID), 1 kb segments & 0.0--0.0\% & solved:1 & 76.8 & 71.0 \\
\bottomrule
\end{tabular}
\par\smallskip
\begin{minipage}{\linewidth}\footnotesize\raggedright
Shared: range over the family's datasets of the fraction of sequences found in GB/NT. Labels: `hard' = MCC below 0.7 with no single route (Table~\ref{tab:gue}). $k$-mer and DNABERT-2: mean MCC$\times$100 (covid: macro-F1$\times$100) of this pipeline's $k$-mer LightGBM and of DNABERT-2 as published; the comparison only locates the baseline. The yeast `negatives' are the measured low-signal class of each mark, not a sampled negative set.
\end{minipage}
\end{table}
